\chapter{Las resistencias y sus límites}

\label{cap:resistencias}
Un libro que sigue expedientes corre el riesgo de describir un territorio donde sólo hay
Estado y mercado. Este capítulo mira lo demás: quiénes reclamaron, qué consiguieron, qué
no, y quiénes no aparecen en ninguna de las fuentes que este trabajo pudo leer. Incluye
además las dos actividades que ocupan el suelo del departamento sin ser ni loteo ni
agricultura —la extracción de áridos y la minería—, porque son las que mejor muestran
sobre qué se monta todo lo anterior.

\section{Las comunidades existen, y el Estado lo certificó}

Este libro declaró, entre sus límites, que no pudo verificar la existencia de comunidades
originarias en el departamento ni la de consulta previa. \textbf{Las dos cosas existen, están
certificadas por organismos del Estado y constan en el expediente del amparo.} Corresponde
corregirlo de frente y en el lugar donde más pesa.

\begin{ficha}{Expediente Nº 0010007-67527/2019-28 --- Fiscalía de Estado de la Provincia de Salta}
El 3 de agosto de 2023, la \textbf{Dra.\ Paula Herrando} solicita los antecedentes
territoriales «a fin de convocar a las comunidades originarias a los talleres participativos
para la elaboración e implementación del \textbf{Plan Integral de Manejo de la Microcuenca del
Río La Caldera}». Eleva el pedido el Coordinador General de Fiscalía de Estado,
Dr.\ Pablo G.\ Buccianti, al Subsecretario de Regularización Territorial y Registro de
Comunidades Indígenas, Dr.\ Ariel Sánchez, el 7 de agosto. El alcance 28 del expediente ---abierto en 2019--- se inicia en la mesa de
entradas de Fiscalía el 9, con carátula de pedido «en carácter urgente», y la Secretaría lo
devuelve con los dos informes el 10: \textbf{un día de trámite}.

\textbf{Informe de la Dirección de Registro de Comunidades Indígenas (10 de agosto de 2023).}
El Dr.\ Sebastián Ramiro González ratifica la presencia en el departamento La Caldera de:
\textbf{Comunidad Originaria Kolla Kondorwaira}, con personería reconocida por Resolución
Ministerial Nº 280/2009, en \textbf{Potrero de Castilla --- Vaqueros}; y \textbf{Comunidad
Indígena Kolla Urkuwasi}, con trámite de personería en el expediente 54-142017/2018-0. \textit{La invitación a la consulta que la Subsecretaría le cursa el 6 de septiembre le asigna otro número de personería: «Expte.\ Nº 224-252732/18». El libro no puede decir cuál es el correcto.}

\textbf{Informe del Equipo Técnico de la Ley 26.160 (misma fecha).} La Ing.\ Rosana Serapio
certifica que \textbf{la zona alta de la cuenca del río La Caldera es territorio de ocupación
actual, tradicional y pública de la Comunidad Originaria Kolla Kondorwaira}, según el
relevamiento centralizado del INAI de 2011, aprobado por Resolución INAI Nº 617/2013.
\end{ficha}

Cinco cosas.

\textbf{Primera: hay dos comunidades, y una tiene personería desde 2009.} No es una presunción
ni una reivindicación en disputa: es un registro estatal con número de resolución. La segunda
tiene expediente de personería en trámite desde 2018.

\textbf{Segunda: el territorio está relevado desde 2011 y aprobado desde 2013.} La Ley 26.160
ordenó relevar la ocupación tradicional de las comunidades; el relevamiento se hizo, se
centralizó en el INAI y se aprobó por resolución. \textbf{Ese relevamiento existe desde antes
de que se aprobara ninguno de los loteos que este libro documenta.}

\textbf{Tercera: la comunidad Kondorwaira está en Potrero de Castilla.} El capítulo \ref{cap:fincas} documenta
que «Potrero de Castilla» es uno de los trece partidos policiales en que la provincia dividió
el departamento en 1909, y que el artículo 6º del plan urbano de 2015 lo sigue enumerando
entre las unidades del municipio. \textbf{La unidad administrativa más antigua del archivo y el territorio comunitario reconocido por el INAI llevan el mismo nombre.} Y el río Potrero de Castilla es, por la Ley 4987, el tramo alto del límite entre La Caldera y Vaqueros (capítulo \ref{cap:cot}): eso explica que el registro sitúe a la comunidad en «Potrero de Castilla --- Vaqueros» mientras el plan urbano de La Caldera enumera ese paraje como propio. El paraje está sobre la línea que separa a los dos municipios.

\textbf{Cuarta, y es lo que más importa: la zona alta de la cuenca es territorio de ocupación
tradicional.} La zona alta es la cabecera del río cuyos brazos el informe de 1935
identificó como el origen del problema, aguas arriba de la bifurcación donde ese informe proponía intervenir (capítulo \ref{cap:defensas}); y
es aguas arriba de las tomas que abastecen a la ciudad de Salta.

\textbf{Y quinta: la convocatoria se hizo para los talleres participativos.} El pedido de
agosto de 2023 tenía por objeto convocar a las comunidades, y el expediente judicial registra
en septiembre de ese año las «constancias de consulta libre, previa e informada» y en
noviembre las constancias del taller participativo del 10 de octubre. \textbf{La consulta se hizo, y se hizo en el marco del Plan
de Manejo.}

\subsection*{La consulta tuvo fecha, mesa y participantes}

El Plan de Manejo fecha la consulta y dice quiénes estuvieron. Es una etapa que el
expediente judicial registra como constancia y que aquí aparece descripta.

\begin{ficha}{Plan de Manejo, páginas 124 a 126 --- consulta libre, previa e informada}
«Esta consulta es el compromiso asumido durante la primera fase realizada por funcionarios
provinciales de la \textbf{Secretaría de Asuntos Indígenas, Fiscalía de Estado, Secretaría de
Minería, Secretaría de Recursos Hídricos, Secretaría de Ambiente y Desarrollo Sustentable y
municipalidad de La Caldera}; el \textbf{23 de septiembre de 2023}, donde se presentó la
línea de base ambiental y social que sería utilizada como base para el presente estudio.»

Modalidad: presentación en diapositivas y mapas interactivos. Resultado: «No se presentaron
objeciones a la información brindada por la consultora, reafirmando su interés y ganas de
participar en todas las instancias de este proyecto».
\end{ficha}

\textbf{Cinco organismos provinciales y el municipio, según el Plan; seis, según el acta}, que registra además al Ministerio de Producción y Desarrollo Sustentable. Las invitaciones del 6 de septiembre anunciaban exactamente los cinco del Plan: Asuntos Indígenas, Fiscalía de Estado, Minería, Recursos Hídricos y Ambiente, más el municipio. \textbf{Una corrección:} la fecha que da el Plan no coincide con la de los documentos de la consulta. La convocatoria, las notificaciones del 6 de septiembre, el acta manuscrita y los dictámenes de Asuntos Indígenas sitúan la reunión el \textbf{12 de septiembre de 2023}, en la Municipalidad de La Caldera (capítulo \ref{cap:amparo}).

Y hay una presencia que conviene anotar: \textbf{la Secretaría de Minería}. Este capítulo
documenta seis concesiones de áridos sobre esta cuenca entre 2014 y 2015, tres de ellas sobre tierra fiscal. Que el organismo minero participe de una consulta sobre el manejo de la
microcuenca no es extraño: es coherente con el peso que la extracción de áridos tiene en el
diagnóstico. \textbf{Pero es la primera vez que este libro lo ve sentado en la misma mesa que
la autoridad hídrica y la ambiental.}

Esa «primera fase» y su línea de base ambiental y social son un documento anterior al Plan,
y este relevamiento no lo tiene.

\subsection*{Las canteras tienen nombre, y el Estado las contó}

En 2022, cumpliendo un mandato de la sentencia, la \textbf{Secretaría de Minería y Energía}
elevó la nómina de concesiones vigentes sobre los ríos Vaqueros, La Caldera, Wierna y
Mojotoro, con el estado de sus informes de impacto ambiental. La nota es del \textbf{14 de
junio de 2022}, va dirigida a la Jefa de Programa Gestión y Policía Minera, y \textbf{no lleva
firma en su propia hoja} (lámina \ref{lam:notamineria}). Este capítulo documentaba seis
concesiones a partir de boletines sueltos; el informe oficial trae \textbf{tres cuadros --- de
2021, 2020 y 2018 --- con veinte concesiones en total}, todas consignadas como vigentes. El
facsímil del primero va en la lámina \ref{lam:iia}.

\needspace{9\baselineskip}
{\sffamily\bfseries\footnotesize Informes de impacto ambiental, cuadro de 2021\par}
\vspace{0.5em}
{\small
\setlength{\LTpre}{0.6em}\setlength{\LTpost}{1.2em}
\begin{longtable}{@{}p{3.0cm}p{1.1cm}p{4.0cm}p{3.4cm}@{}}
\toprule
\textbf{Cantera} & \textbf{Expte.} & \textbf{Concesionario} & \textbf{Estado del IIA} \\
\midrule
\endhead
La Mesa Redonda & 14971 & Cooperativa Mesa Redonda & \textbf{No aprobado} \\
Tierra Gaucha & 20738 & Geomix S.R.L. & P/firma Secretario \\
Mojotoro Norte & 22776 & La Población de Ortega S.A. & P/firma Secretario \\
Don Roque I & 22290 & MEI Obras y Servicios S.R.L. & \textbf{Aprobado, Res.\ 119} \\
La Serena & 18995 & Alejandrina Morales & P/firma Secretario \\
La Mesa Redonda I & 18134 & Cooperativa Mesa Redonda & P/firma Secretario \\
Ana María & 17033 & Carlos Fernández Pérez & P/firma Secretario \\
Puente Ferroviario & 16386 & Cooperativa Mesa Redonda & P/firma Secretario \\
La Ferroviaria & 17394 & Cooperativa Mesa Redonda & P/firma Secretario \\
El Triunfo II & 17422 & Cooperativa Mesa Redonda & P/firma Secretario \\
Triunfo III & 17532 & Cooperativa Mesa Redonda & P/firma Secretario \\
La Generosa & 16893 & Héctor Enrique Medina & \textbf{No aprobado} \\
Tomasito & 17049 & Héctor Enrique Medina & \textbf{No aprobado} \\
Teresa & 17508 & Gonzalo Federico Ventícola & \textbf{Aprobado, Res.\ 099} \\
\bottomrule
\end{longtable}
}

\needspace{9\baselineskip}
{\sffamily\bfseries\footnotesize Cuadros de 2020 y de 2018\par}
\vspace{0.5em}
{\small
\setlength{\LTpre}{0.6em}\setlength{\LTpost}{1.2em}
\begin{longtable}{@{}p{3.0cm}p{1.1cm}p{4.0cm}p{3.4cm}@{}}
\toprule
\textbf{Cantera} & \textbf{Expte.} & \textbf{Concesionario} & \textbf{Estado del IIA} \\
\midrule
\endhead
La Cornisa & 17088 & Juan Carlos Lozano & \textbf{Aprobado, Res.\ 048} \\
Noroeste Construcciones & 19065 & Noroeste Construcciones S.A. & \textbf{No aprobado} \\
La Vaquera & 18770 & César Domingo Dal Borgo & \textbf{No aprobado} \\
Servidumbre de campamento, acopio y clasificación & 21615 & César Domingo Dal Borgo & \textbf{No aprobado} \\
\midrule
Macarena & 18219 & Hermanos Mendoza S.R.L.\textsuperscript{a} & \textbf{Aprobado, Res.\ 276} \\
Los Yacones & 20604 & INCOVI S.R.L. & \textbf{Aprobado, Res.\ 340} \\
\bottomrule
\end{longtable}
}

\begin{cautela}
La transcripción se hizo sobre el facsímil reproducido en la lámina \ref{lam:iia} y sobre la
hoja siguiente del mismo informe, que trae los cuadros de 2020 y 2018. \textbf{Una versión
anterior de este apartado transcribía solo las ocho primeras filas de 2021 y advertía que las
restantes no se habían verificado.} Ya están, y con ellas los otros dos cuadros. Hay, además, una discordancia entre fuentes: Los Yacones, aprobada y vigente en estos cuadros, figura en el informe de Minería de 2023 con la renovación rechazada en junio de 2020; y La Mesa Redonda I, «para firma del Secretario» en el cuadro de 2021, figura en ese mismo informe con el informe de impacto aprobado (capítulo \ref{cap:amparo}), lo que puede indicar una aprobación posterior o un error en una de las dos fuentes. Se mantiene
el criterio: el cuadro debe cotejarse contra el original antes de citar cualquier fila ---\textsuperscript{a} los edictos de 2014 nombran a la concesionaria «Transporte Hnos.\ Mendoza S.R.L.»; que sea la misma sociedad es probable y no está verificado---, y las
columnas de fecha de entrada y de municipio del documento no se reproducen acá por razones de
composición. \textbf{Del reparto por municipio, que el apartado usa más abajo, se consigna el agregado}: dieciséis de las veinte filas están asignadas a «La Caldera-Vaqueros» y las cuatro restantes a «La Caldera».
\end{cautela}

\textbf{Primera: de las catorce concesiones de 2021, dos tienen el informe de impacto ambiental
aprobado.} Nueve están «para firma del Secretario» --- una desde marzo de 2020 --- y tres
figuran expresamente como \textbf{no aprobadas}. Las catorce constan como \textbf{vigentes}. En
el cuadro de 2020 la proporción es peor: de cuatro, tres no aprobadas, y las cuatro vigentes.
\textbf{Sobre veinte concesiones, cinco tienen la evaluación ambiental aprobada y seis constan
como no aprobadas; las veinte están vigentes.}

\textbf{Segunda: una concentración que el cuadro deja ver.} La \textbf{Cooperativa Mesa Redonda}
figura como concesionaria en \textbf{seis} de las catorce filas de 2021 --- La Mesa Redonda, La
Mesa Redonda I, Puente Ferroviario, La Ferroviaria, El Triunfo II y Triunfo III ---, y ninguno
de esos seis informes está aprobado: cinco esperan firma y uno consta como no aprobado. César
Domingo Dal Borgo aparece dos veces en 2020 y Héctor Enrique Medina dos veces en 2021, en los
cuatro casos con el informe no aprobado. El capítulo documenta más abajo, por boletines
sueltos, la renovación de «El Triunfo II» a nombre de la Cooperativa de Provisión de Áridos La
Mesa Redonda Ltda.\ bajo el expediente 17.422: \textbf{es la misma fila del cuadro oficial}, y
las dos fuentes se confirman en el número, el nombre y el concesionario. \textbf{No se confirman
en la jurisdicción}: el edicto de abril de 2014 ubica esa cantera «en el departamento Capital,
río La Caldera», mientras la columna de municipio del cuadro sólo admite «La Caldera» y «La
Caldera-Vaqueros». Una de las dos fuentes le asigna al departamento una cantera que la otra
ubica fuera de él. No es un detalle ocioso: el cuadro es la base del recuento de veinte
concesiones que este libro cita en otros cinco capítulos. \textbf{Ese número debe leerse, de aquí en adelante, con esta reserva}: son veinte según el registro del propio organismo, y al menos una de ellas ---«El Triunfo II»--- aparece en el Boletín Oficial ubicada en el departamento Capital. Cuántas más estén en la misma situación es algo que el cuadro no permite saber, porque su columna de municipio sólo admite dos valores y ninguno de los dos es «Capital».

\subsection*{Y el edicto minero sí publica las coordenadas}

\begin{ficha}{Boletín Oficial, 17 y 26 de mayo y 7 de junio de 2016 --- Expte.\ Nº 14.971}
El Juez de Minas «hace saber a los efectos del \textbf{art.\ 73 del C.P.M.\ Ley 7141/01} que la \textbf{Cooperativa de Provisión Minera «La Mesa Redonda» Ltda.} ha solicitado publicación de la \textbf{renovación de concesión} de la cantera de áridos denominada \textbf{La Mesa Redonda} \ldots\ ubicada en el departamento La Caldera, \textbf{lugar Río La Caldera}».

Sigue la descripción: \textbf{cuarenta y cinco pares de coordenadas Gauss Krüger}, con la aclaración «Sistema Posgar 94 y Campo Inchauspe/69».

\textbf{Superficie libre: 32 hectáreas 5.240 metros cuadrados. «Los terrenos afectados son de propiedad fiscal.»}

Firma la Dra.\ M.\ Alejandra Loutayf, secretaria. Importe: \$300.
\end{ficha}

\textbf{Ese edicto contesta dos cosas que el cuadro no decía.} La cantera «La Mesa Redonda» del expediente 14.971 ---la que el cuadro de 2021 registra con informe \textbf{no aprobado}--- ocupa \textbf{treinta y dos hectáreas y media de tierra fiscal sobre el río La Caldera}.

\textbf{Y corrige, por acumulación, una cifra que este libro venía repitiendo.} La fórmula que viaja a otros capítulos ---«seis concesiones entre 2014 y 2015, tres de ellas sobre tierra fiscal»--- es correcta para esa ventana y deja afuera, por un año, la mayor de todas. Puestas juntas, las canteras sobre tierra pública que este relevamiento documenta son \textbf{cuatro} y suman \textbf{setenta y tres hectáreas}: «Macarena» (10 ha), «El Vaquero» (9 ha), «Los Yacones» (21,6 ha) y «La Mesa Redonda» (32,5 ha). \textbf{La más grande es la que tiene el informe de impacto expresamente no aprobado.}

\begin{cautela}
\textbf{Y trae una comparación que este libro debe hacer, porque la tiene servida.}

El Juzgado de Minas publica \textbf{cuarenta y cinco pares de coordenadas} para delimitar una cantera, por trescientos pesos, en 2016. La autoridad hídrica, sobre los mismos cauces y en los mismos años, publica once pares en 2013, ciento trece en 2014 y \textbf{ninguno} en diciembre de ese mismo año \textit{(cap.\ \ref{cap:ribera})}.

\textbf{Dos organismos de la misma provincia, sobre el mismo río, con criterios distintos de publicidad.} Y el que delimita para extraer publica más que el que delimita para proteger.
\end{cautela}

\subsection*{Octubre de 2018: la gacetilla vecinal}

Antes de seguir con los cuadros, corresponde registrar que esta actividad tuvo oposición
organizada, y que la tuvo mucho antes de que este libro existiera.

\begin{ficha}{Manifiesto de Vecinos Autoconvocados de La Caldera, La Calderilla y Vaqueros, difundido en octubre de 2018}
El documento denuncia la \textbf{extracción de áridos «en toda la zona de ribera desde el puente
Wierna hasta el mismo pueblo de La Caldera»} y señala que el reclamo \textbf{lleva cinco años}
sin respuesta del municipio, del ministerio ambiental ni de la autoridad minera ---«un reclamo que cumple ya 5 años», en el texto publicado---.

Sobre la convivencia entre extracción y vivienda: «\textit{Existen actualmente numerosas
urbanizaciones que colindan con el rio La Caldera en áreas de extracción de áridos, dándose el
absurdo de tener \textbf{maquinaria pesada trabajando a pocos metros de casas en barrios
como el Duranzo} [sic] \textbf{y el Nogalar}, así como también a lo largo de toda la costa ribereña}».

Enumera como consecuencias la contaminación sonora, la degradación del paisaje y \textbf{«el
riesgo de inundación en temporada de lluvias o crecida del caudal, producto de la degradación de
las barreras naturales del cauce del rio»}.

Y describe el reparto de competencias: las áreas, los volúmenes y las reglas las define la
autoridad minera provincial, \textbf{que puede convenir con los municipios el control y el cobro
de un canon}; «\textit{aunque este convenio existe, el reclamo gira en torno al control efectivo
que debería realizar el Municipio}».

\textit{Fuentes: \textit{Cuarto Poder}, «Denuncian extracción desmedida de áridos en La Caldera, Calderilla y Vaqueros», 16 de octubre de 2018, \url{https://cuartopodersalta.com.ar/denuncian-extraccion-desmedida-de-aridos-en-la-caldera-calderilla-y-vaqueros/}, de donde salen las citas, cotejadas letra por letra el 18 de septiembre de 2026 ---«rio» y «Duranzo» están así en el original---; y \textit{La Gaceta Salta}, «Vecinos de La Caldera repudian la extracción de áridos», 18 de octubre de 2018, \url{https://www.lagacetasalta.com.ar/nota/112667/actualidad/vecinos-caldera-repudian-extraccion-aridos.html}, que atribuye a los vecinos, entre comillas, el pasaje sobre la maquinaria. Cuarto Poder publica el texto sin comillas y bajo su firma, de modo que los otros pasajes pueden ser del manifiesto o del medio; el libro los atribuye al manifiesto porque la otra nota confirma su origen para el primero. El título «Por ríos limpios, seguros, vivos y de todos», con el que una versión anterior de esta ficha encabezaba el documento, no figura en ninguna de las dos notas.} \pendiente{texto original del manifiesto vecinal de octubre de 2018, con su título y la nómina de firmantes}
\end{ficha}

\textbf{Tres cosas de esa gacetilla son verificables contra lo que este libro ya documenta, y
las tres coinciden.}

\textbf{El tramo.} «Desde el puente Wierna hasta el pueblo» es exactamente donde el cuadro
oficial ubica las concesiones de esta cooperativa \textit{(supra)}.

\textbf{Los barrios.} El Durazno y El Nogalar son dos de los loteos actores del amparo; el primero es el que el capítulo \ref{cap:loteo} documenta con expediente, y sus arroyos son los que determinó sin publicar coordenadas la Resolución 383/14 \textit{(cap.\ \ref{cap:ribera})}.

\textbf{Y el riesgo.} «La degradación de las barreras naturales del cauce» en temporada de
crecida es, dicho en 2018 y por vecinos, lo mismo que el capítulo \ref{cap:defensas} documenta
con expedientes de defensas y el capítulo \ref{cap:amparo} con un amparo colectivo.

\begin{cautela}
\textbf{La gacetilla contiene además una imputación personal, y este libro la trata como lo que
es.}

Los vecinos denunciaron, según las dos notas que la difunden, que el entonces intendente
del municipio \textbf{integraba la cooperativa que debía controlar}.

\textbf{Este libro no verificó esa afirmación y no la hace suya.} No tiene la nómina de
asociados de la cooperativa, no tiene constancia registral alguna, y \textbf{no dio al
interesado la oportunidad de responder}. Lo que consta es que \textbf{la denuncia se formuló,
que dos medios la publicaron en octubre de 2018 y que este relevamiento la encontró}: nada más.

\textbf{Lo que la resolvería es un solo documento}: la nómina de asociados de la cooperativa
\textbf{matrícula 22.359} a la fecha de la denuncia, en el registro de cooperativas. Está en el
apéndice \ref{cap:pedidos}.

\textbf{Y si ese documento apareciera}, el capítulo \ref{cap:politica} fija lo que corresponde
antes de publicar nada: la carta de derecho a réplica, enviada y con plazo cumplido. \textbf{Una
denuncia vecinal no releva de ese paso; lo hace más necesario.}
\end{cautela}

\begin{ficha}{Boletín Oficial Nº 21.447, 11 de abril de 2023 --- identificación de la cooperativa}
«La Comisión Directiva de la \textbf{Cooperativa de Provisión Minera “La Mesa Redonda Ltda.”},
\textbf{Matrícula 22.359}, convoca a los señores asociados a la Asamblea General Ordinaria
\ldots\ en la \textbf{sede social de Zuviría Nº 2370}», para tratar el \textbf{ejercicio
económico Nº 22, cerrado el 31 de diciembre de 2022}.

Firman \textbf{Juan Domingo Lozano}, presidente, y \textbf{Juan José Vargas}, secretario.
\end{ficha}

\textbf{Eso agrega tres cosas al cuadro.} La \textbf{matrícula 22.359}, que la vuelve pedible en
el registro de cooperativas. El \textbf{domicilio}, en la ciudad de Salta. Y el número de
ejercicio: veintidós cerrados en 2022 la ubica \textbf{constituida hacia el año 2000},
\textbf{anterior a todas las concesiones que el cuadro le asigna}.

\begin{cautela}
\textbf{Y una discrepancia de denominación que conviene registrar.} El boletín de la renovación
de «El Triunfo II» la nombra \textbf{Cooperativa de Provisión de Áridos La Mesa Redonda Ltda.};
el de la asamblea de 2023, \textbf{Cooperativa de Provisión Minera «La Mesa Redonda» Ltda.}

Puede ser la misma entidad nombrada de dos modos, un cambio de denominación social, o dos
cooperativas distintas. \textbf{Este libro no lo establece}, y la matrícula 22.359 es lo que
permite dirimirlo en un solo trámite.
\end{cautela}

\textbf{Tercera: dieciséis de las veinte están asignadas a «La Caldera-Vaqueros».} Los cuadros
no tienen columna de río: tienen columna de \textbf{municipio}, y solo admite dos valores ---
«La Caldera» y «La Caldera-Vaqueros». Más de tres cuartos de la actividad extractiva registrada
recae, según el propio organismo, sobre \textbf{una jurisdicción compartida entre los dos
municipios}. Es una frontera del mismo tipo que la que el capítulo \ref{cap:vaqueros} documenta como tierra de nadie administrativa ---la de los ríos Wierna y La Caldera, que el capítulo \ref{cap:cot} reconstruye desde las leyes---, y ahora se sabe qué se extrae de ella.

\begin{cautela}
Hay una discordancia entre la nota y sus propios cuadros que conviene no pasar por alto. El
punto 1 anuncia las concesiones «sobre los ríos Vaqueros, La Caldera, Wierna, Mojotoro». Los
cuadros \textbf{no traen ninguna columna de río}: ninguna fila queda asignada al Wierna ni al
Mojotoro como curso. La única aparición del Mojotoro es el \textbf{nombre} de una cantera
---«Mojotoro Norte»---, y su municipio es La Caldera. \textbf{Lo que la nota promete por río,
el anexo lo entrega por municipio}, y los dos recortes no son traducibles uno al otro.
\end{cautela}

\textbf{Cuarta: la concesión la otorga un organismo y la evaluación ambiental la hace otro.}
El informe lo dice: «La concesión de las canteras las realiza el Juzgado de Minas de la
Provincia», mientras que los informes de impacto «son evaluados por personal de este
Programa». \textbf{La cantera puede estar vigente con su evaluación ambiental pendiente de
firma, porque quien concede no es quien evalúa.}

\textbf{Y quinta, que es la más grave y la dice el propio Estado:}

\begin{ficha}{Secretaría de Minería y Energía, nota del 14 de junio de 2022 --- punto 6}
«Se hace notar que en estos ríos se observa \textbf{gran cantidad de loteos y asentamientos de
viviendas en zonas inundables y dentro de las líneas de ribera delimitada por la Secretaría de
Recursos Hídricos}. Situación que debe ser atendida de manera urgente por la problemática que
ocasiona.»
\end{ficha}

\textbf{Y lo dice un organismo sobre el trabajo de otro.} No es el mismo que traza la línea:
es la autoridad minera informando, en una pieza que terminó en un expediente judicial, que la
línea trazada por la autoridad hídrica está ocupada por loteos y asentamientos. No es una
inferencia de este libro ni un reclamo vecinal, y \textbf{no puede atribuirse a un organismo
interesado en defender su propia delimitación}, porque el que lo dice no la hizo.

Y hay un detalle técnico que conviene retener. La nota precisa que la normativa exige que la
extracción se haga «por raspado a profundidades \textbf{hasta el metro}» y que el descarte se
deposite «a modo de defensa sobre ambas márgenes». \textbf{La presentación de la Secretaría de
Recursos Hídricos ante el Consejo Hídrico Federal dice que en algunos casos se extrae por
debajo de un metro y medio.} Un organismo fija el metro; el otro informa que se extrae por
debajo de un metro y medio.

\begin{cautela}
\textbf{Y hay un punto 3 en esa misma nota que conviene leer contra los cuadros.} Dice que
dentro de la \textbf{Zona de Exclusión de Aguas del Norte} existía un convenio con esa empresa
para ordenar los laboreos, que \textbf{a la fecha ese convenio no se encuentra vigente} y que
por esa razón \textbf{se paralizaron las actividades extractivas}, girándose las actuaciones a
la Dirección General de Asuntos Legales y Técnicos.

En la misma pieza, las veinte concesiones de los tres cuadros figuran con estado
\textbf{«Vigente»}. Este libro no puede establecer si la parálisis alcanza a todas, a algunas o
solo al tramo de la zona de exclusión: \textbf{la nota no lo dice y los cuadros no lo
distinguen}. Se consigna la convivencia de las dos declaraciones sin conciliarlas.
\end{cautela}

\textbf{Una última cosa sobre quién firma.} Los tres bloques de firma del anexo --- el Ing.\
Guillermo A.\ Fuchs, el Ing.\ Civil Ricardo Cristóbal Ramos y Mauricio Romero Leal --- son de
la \textbf{Secretaría de Recursos Hídricos}, no de la Secretaría de Minería y Energía, que es
la que emite la nota y a la que pertenecen los datos. Son los mismos tres que firman el informe
de mapas del capítulo \ref{cap:amparo}. \textbf{El anexo minero está firmado por funcionarios de otra
secretaría}, y la nota que lo encabeza no está firmada por nadie. Se registra el hecho; este
libro no lo explica.

\subsection*{Y el Estado las puso en un mapa}

Las canteras no solo tienen nombre: tienen ubicación oficial. El mismo informe de mayo de 2022
trae un tercer mapa temático --- \textbf{lámina \ref{lam:mapa3}} --- con los polígonos mineros
sobre las márgenes del río La Caldera y sus cauces aportantes, a escala 1:40.000 y en POSGAR
2007, Faja 3. Se armó, dice el informe, sobre la capa «Catastro Minero: Polígonos», que es la
representación geográfica de los datos contenidos en los expedientes judiciales, y distingue
los polígonos de cantera de los de convenio.

Sobre el plano se leen nombres que este capítulo viene siguiendo: \textbf{Tierra Gaucha} y
Tierra Gaucha 20738 al norte, \textbf{Mesa Redonda} y \textbf{La Mesa Redonda I} sobre el
tramo medio, y hacia el sur del plano, en el sector del Wierna próximo a su unión con el río La Caldera, \textbf{Los Yacones}, \textbf{Terranostra},
\textbf{Hanna} y \textbf{Noroeste Construcciones}. Figura además un polígono rotulado
\textbf{Aguas del Norte S.A.}

\begin{cautela}
\textbf{La lámina no permite medir superficies ni verificar linderos} a la resolución del
ejemplar incorporado al expediente, y los rótulos del plano no coinciden con las filas de
los cuadros de informes de impacto. \textbf{Dos polígonos rotulados en el mapa ---
«Terranostra» y «Hanna» --- no figuran en ninguno de los tres cuadros.} Las dos nóminas
provienen de fuentes distintas --- el catastro minero una, la evaluación ambiental la otra ---
y el cotejo entre ambas requiere las capas vectoriales. El informe de Minería de 2023 la explica en parte ---Hanna con concesión provisoria y Terranostra con la concesión vencida (capítulo \ref{cap:amparo})---; para el resto, se consigna la discrepancia sin
explicarla: puede tratarse de polígonos con expediente y sin informe presentado, de convenios
en lugar de canteras, o de registros de años no cubiertos por los cuadros.

Lo que sí queda establecido, y no es poco, es que \textbf{la extracción de áridos de esta
cuenca está cartografiada por el Estado desde 2022, en el mismo sistema de coordenadas y a la
misma escala que las líneas de ribera} de la lámina \ref{lam:mapa2}. Las dos capas son
superponibles. Que lo sean vuelve verificable, punto por punto, la afirmación del punto 6 de
la nota transcripta más arriba: que hay loteos y asentamientos dentro de las líneas de ribera
que el propio organismo delimitó.
\end{cautela}

\lamina{lam-mapa3-mineros.jpg}{lam:mapa3}%
{Polígonos mineros sobre el río La Caldera y sus cauces aportantes, 2022}%
{\textbf{Mapa Nº 3.} Ubicación de los polígonos mineros en márgenes del río La Caldera y sus
cauces aportantes, aguas arriba de la unión con el río Wierna, registrados en la Secretaría de
Minería hasta el año 2022. Escala 1:40.000, sistema POSGAR 2007 --- Faja 3. Distingue los
polígonos de cantera de los de convenio, sobre la capa «Catastro Minero: Polígonos».}%
{«Informe de mapas de cuenca de río La Caldera», Secretaría de Recursos Hídricos de Salta, mayo
de 2022, presentado en la ejecución de sentencia del amparo colectivo, expediente
INC-23860/1. Reproducción facsimilar; se recortó el margen en blanco del escaneo y no se
alteraron los tonos.}

\lamina{lam-nota-mineria.jpg}{lam:notamineria}%
{Nota de la Secretaría de Minería y Energía al juzgado, 14 de junio de 2022}%
{\textbf{Nota del 14 de junio de 2022}, dirigida a la Jefa de Programa Gestión y Policía Minera, que acompaña los tres cuadros de informes de impacto ambiental. Su punto 6 hace constar «gran cantidad de loteos y asentamientos de viviendas en zonas inundables y dentro de las líneas de ribera delimitada por la Secretaría de Recursos Hídricos». \textbf{La hoja no lleva firma}, y el anexo que la acompaña está firmado por funcionarios de otra secretaría.}%
{Ministerio de Producción y Desarrollo Sustentable, Secretaría de Minería y Energía, nota incorporada al expediente INC-23860/1. Reproducción facsimilar; se recortó el margen en blanco del escaneo y no se alteraron los tonos.}

\lamina{lam-iia-2021.jpg}{lam:iia}%
{Cuadro de informes de impacto ambiental de canteras, 2021}%
{\textbf{Cuadro «IIA 2021»} de la Secretaría de Minería y Energía: catorce concesiones de
áridos en La Caldera y Vaqueros, con expediente, concesionario, resultado del informe de
impacto ambiental, fecha de entrada y estado. Nueve figuran «p/firma Secretario», tres como no
aprobadas, dos aprobadas por resolución. Las catorce, vigentes.}%
{Ministerio de Producción y Desarrollo Sustentable, Secretaría de Minería y Energía, anexo a la
nota del 14 de junio de 2022, expediente INC-23860/1. Reproducción facsimilar; se recortó el
margen en blanco del escaneo y no se alteraron los tonos.}

\subsection*{La extracción de áridos, descripta por el organismo hídrico}

Este capítulo documenta seis concesiones de áridos sobre la cuenca entre 2014 y 2015, y la
prensa atribuye a esa actividad el origen del amparo. La Secretaría de Recursos Hídricos
describió el problema ante el Consejo Hídrico Federal en estos términos:

\begin{ficha}{Torres, J.\ G., Secretaría de Recursos Hídricos de Salta}
«Otros de los problemas que se plantea en los ríos de la provincia, es la extracción de
áridos. \textbf{No por la acción misma, sino por el mal manejo de las operaciones propias de
la actividad y el depósito de los descartes}.»

«Si bien, el Juzgado de Minas y la Secretaria de Minería, están exigiendo a los concesionarios
de las canteras, la delimitación de la línea de ribera \ldots\ \textbf{No cumplen, por lo
general, con las normativas vigentes}.»

«La extracción, en algunos casos, se realiza \textbf{por debajo de 1,5 m de profundidad}. Esto
trae aparejado una modificación de los canales de circulación de agua del sub-álveo.»

«La acumulación de descartes y playa de maniobras, \textbf{se realizan dentro de la línea de
ribera} y ocasionan una \textbf{disminución del ancho del cauce}, y trae aparejado problemas
de desbordes.»
\end{ficha}

Tres cosas quedan dichas por el Estado que este libro solo podía insinuar.

\textbf{Que los concesionarios, por lo general, no cumplen.} No es una denuncia vecinal ni una
conclusión de este trabajo: es la evaluación del organismo hídrico sobre la actividad que
otro organismo concesiona.

\textbf{Que el incumplimiento produce desbordes.} Los descartes dentro de la línea de ribera
reducen el ancho del cauce. \textbf{Es la cadena causal completa entre la extracción y la
inundación}, enunciada por la Secretaría de Recursos Hídricos.

\textbf{Y que la exigencia existe.} El Juzgado de Minas y la Secretaría de Minería piden la
delimitación de línea de ribera a los concesionarios de canteras. El capítulo \ref{cap:ribera} documenta
veinticinco actos de línea de ribera sobre esta cuenca entre 2013 y 2015; \textbf{ahora se entiende que una parte pudo motivarla la extracción}, aunque la mayoría de esos actos recae sobre fincas y loteos.

\subsection*{La comunidad publicó un libro, y nombra al río}

En 2019 la Comunidad Kolla Kondorwaira publicó \textit{Nuestro tiempo redondo}, un libro de
cincuenta y siete páginas sobre su territorio y su calendario, editado por el propio
municipio de La Caldera. Es la única pieza de este relevamiento \textbf{escrita por habitantes
del departamento sobre el departamento}, y por eso merece leerse con cuidado.

\subsubsection*{Primero: cómo se define a sí misma}

\begin{ficha}{\textit{Nuestro tiempo redondo}, sección «La Comunidad Kondorwaira»}
La comunidad se sitúa en \textbf{Potrero de Castilla, departamento La Caldera}, en «la alta
cuenca del \textbf{río Mojotoro}», de bosques montanos a pastizales de altura, «a escasos
kilómetros de Salta Capital».

Y define su propia dinámica territorial: incorpora «las localidades urbanas de
\textbf{Vaqueros y La Caldera} y barrios de la \textbf{zona norte de la ciudad de Salta}».

\textit{Kondorwaira} es quechua: \textbf{cóndor al viento}.

En la sección metodológica, firmada por Karla Pérez Domínguez y Silvina Belmonte, se precisa
que la comunidad habita «la zona alta de la cuenca del \textbf{Río Las Nieves}», en un
territorio «de base agrícola-ganadera, con manejo de bienes comunes en temporalidad cíclica».
\end{ficha}

\textbf{La primera consecuencia es que la comunidad describe el mismo corredor que este libro,
y lo describe como una unidad.} Vaqueros, La Caldera y el norte de la ciudad no son para ella
tres jurisdicciones: son el recorrido de sus bajadas. \textbf{Es la definición funcional del
corredor norte, enunciada por quienes lo caminan}, y no aparece en ninguno de los instrumentos que los capítulos \ref{cap:pdua} y \ref{cap:cot} analizan; el único acto que trató a esos municipios como una unidad es el convenio del Corredor Intermunicipal de 2009 (capítulo \ref{cap:redes}), del que no se conocen actos de funcionamiento.

\textbf{Y la segunda toca el capítulo \ref{cap:ribera} de lleno: el río Las Nieves.} Es el curso sobre el que
funcionó, de 1947 a 1961, la estación de aforo «El Volcán», con trescientos kilómetros
cuadrados de área de aporte. \textbf{La comunidad nombra como cuenca de habitación el mismo
río que el Estado dejó de medir hace sesenta y cinco años.}

\subsubsection*{Y después: los ríos ordenan el calendario, y no tienen nombre}

El libro menciona ríos en cuatro pasajes y \textbf{en ninguno los nombra}. No hace falta: son
«los ríos», y lo que importa de ellos es cuándo crecen.

\begin{cautela}
Se vuelve del puesto «antes de que crezcan los ríos». Se eligió septiembre --- y después
octubre --- para devolver la Virgen porque «todavía no están crecidos los ríos». Cuando
llueve, no se puede salir.

\textbf{Ese es el sistema de alerta temprana que existe en esta cuenca, y es un calendario.}
El Plan de Manejo propone cinco estaciones hidrométricas para dar cuarenta y cinco minutos de
aviso. La comunidad organiza su año para no estar del lado equivocado del río cuando crezca.
\textbf{Son dos formas de la misma previsión, separadas por un siglo de asimetría técnica.}
\end{cautela}

El libro registra además el \textbf{despoblamiento del cerro y la migración de los jóvenes a la
ciudad}. El capítulo \ref{cap:poblacion} mide ese movimiento desde los censos, como saldo; \textbf{aquí está
dicho desde el lugar del que se van}.

\begin{cautela}
Este libro no usó \textit{Nuestro tiempo redondo} como fuente de datos catastrales porque no
los tiene: \textbf{no hay decretos, expedientes, matrículas, superficies ni importes}, y no
corresponde exigirle lo que no se propuso.

Tampoco resuelve la pregunta abierta de este capítulo --- dónde está el perímetro relevado ---
ni contiene coordenadas.

Se incorpora por lo que sí hace: \textbf{nombrar el territorio desde adentro}. Y conviene
anotar una discrepancia de escala sin forzarla: el libro dice «cuenca del Mojotoro» en un
lugar y «cuenca del Río Las Nieves» en otro. \textbf{No es contradicción --- el Las Nieves
tributa al sistema del Mojotoro --- pero el libro no lo explicita, y la conciliación es lectura
de este trabajo, no del documento.}

Y queda una tercera denominación, que es la del certificado estatal: el Equipo Técnico de la Ley 26.160 sitúa el territorio en «la zona alta de la cuenca del \textbf{río La Caldera}». Las tres nombran el mismo sistema ---el Las Nieves desagua en el Wierna y el Wierna en el río La Caldera, según las leyes de límites que transcribe el capítulo \ref{cap:cot}---, pero ninguno de los documentos lo dice. \textbf{La cadena es, otra vez, lectura de este trabajo.}
\end{cautela}

\subsection*{Veintisiete familias en los cerros}

Este libro nombró a la Kondorwaira como un dato del expediente: una personería, un
relevamiento, una resolución. Un parte del Gobierno provincial de octubre de 2017 permite
decir quiénes son.

\begin{ficha}{Secretaría de Comunicación, Gobierno de Salta, 31 de octubre de 2017}
«La comunidad Kolla Kondorwaira \ldots\ está situada en el paraje \textbf{Potrero de Castilla},
departamento La Caldera, y la conforman \textbf{27 familias que residen en los cerros}.»

Su cacique es \textbf{Teodora Sarapura}, que recibió un segundo aporte no reintegrable de
\textbf{\$1.028.000} del \textbf{Programa de Desarrollo Rural Incluyente (PRODERI)}.

Destino de los fondos: cerramiento de corrales y refugios para \textbf{ovinos y caprinos},
equipos para \textbf{producción de quesos}, compra de reproductores y refacciones. Con el
primer aporte habían comprado ganado menor, herramientas, mejorado pasturas y adquirido
\textbf{una rueca y telares} para trabajar la lana.

Asistencia técnica: Nación, Provincia e \textbf{INTA}, con el coordinador de la Secretaría de
Agricultura Familiar Ricardo Bima.
\end{ficha}

Tres cosas.

\textbf{Primera: son veintisiete familias ---treinta y una según un trabajo académico de 2024 (véase más abajo)---.} A las 3,04 personas por hogar que el Censo 2022 da para esa fracción, \textbf{unas ochenta personas}. \textbf{Es una fracción pequeña de los doce mil
habitantes del departamento, y ocupa la zona alta de la cuenca} --- justamente el sector que
el Plan de Manejo identifica como área de aporte.

\textbf{Segunda: su economía es ganadería menor y lana.} Cabras, ovejas, quesos, una rueca,
telares. Es el único uso productivo que este libro documenta en la cabecera de la cuenca, y no
aparece en ninguno de los instrumentos de ordenamiento territorial que analizan los capítulos \ref{cap:pdua} y \ref{cap:cot}.

\textbf{Y tercera, que es la que más pesa: el Estado provincial financió a esta comunidad al
menos dos veces.} No se trata de un actor que la administración desconociera. \textbf{El
Ministerio de Ambiente y Producción Sustentable le entregó cheques en actos públicos en 2017};
el Ministerio de Asuntos Indígenas le entregó su carpeta técnica en 2016. \textbf{Y en agosto
de 2023 la abogada de Fiscalía de Estado tuvo que pedir «en carácter urgente» que un ministerio
le informara a otro si había comunidades en la cuenca.}

\begin{cautela}
Los datos de este apartado provienen de partes de prensa oficiales, no de instrumentos
públicos. El número de familias, el monto y el destino de los fondos son los que el Gobierno
provincial comunicó en 2017; este libro no accedió al expediente del PRODERI ni verificó la
ejecución.

Tampoco corresponde leer esta sección como un reproche a la comunidad ni como un argumento
sobre sus derechos. \textbf{Se incorpora porque muestra que la información existía, dispersa
entre organismos que no la cruzaron.}
\end{cautela}

\subsection*{Y hay un mapa nacional, con su límite}

El INAI publica tres mapas del país elaborados con datos del Registro Nacional de Comunidades
Indígenas y del programa de relevamiento territorial. Uno de ellos, el \textbf{Mapa de
Territorios}, «visualiza los territorios de ocupación actual, tradicional y pública de las
Comunidades Indígenas que han sido relevados». Están en formato PDF, actualizados a septiembre
de 2024 y son de descarga libre.

\textbf{Sirven para confirmar que el territorio relevado figura en la cartografía oficial de
la Nación. No sirven para la pregunta de este libro.} Son mapas de escala nacional --- seis
megabytes para todo el país --- y a esa escala un departamento entero ocupa unos milímetros:
no permiten establecer si una parcela determinada cae dentro o fuera de un perímetro.

\textbf{Para eso hace falta la carta de la carpeta técnica}, que es lo que el apartado
siguiente documenta.

\subsection*{La carpeta se entregó en 2016, y la comunidad la tiene}

Un parte de prensa del Gobierno de Salta del 25 de noviembre de 2016 documenta el acto en que
se entregaron las carpetas técnicas de relevamiento territorial a cuatro comunidades kollas.
Una de ellas es la de este departamento.

\begin{ficha}{Secretaría de Comunicación, Gobierno de Salta, 25 de noviembre de 2016}
El secretario de Asuntos Indígenas, \textbf{Enrique Rojo}, encabezó la entrega de cuatro
carpetas técnicas a comunidades kollas de los departamentos de Orán, Santa Victoria y
\textbf{La Caldera}: Finca Las Juntas, Finca San Ignacio, El Arazal y \textbf{Asociación
Comunitaria Kondorwaira}.

«El relevamiento constituye una importante herramienta del Estado para identificar las tierras
que tradicionalmente ocupan las comunidades.»

Participaron el subsecretario \textbf{Ariel Sánchez} --- el mismo funcionario a quien Fiscalía
de Estado dirigirá su pedido siete años después --- y representantes del INAI.

Se aclara que en Salta el relevamiento se ejecuta «a través del \textbf{Equipo Técnico
Operativo} conformado por técnicos propuestos por la provincia, el INAI y representantes
indígenas».
\end{ficha}

Tres cosas, y la primera corrige a este libro.

\textbf{Primera: la modalidad no está establecida, y el propio expediente se contradice.} El informe del 10 de agosto de 2023 de la coordinadora provincial del programa dice que el relevamiento «fue realizado por el equipo técnico del INAI, en la \textbf{modalidad centralizado} el día 20 de noviembre del 2011», y define ese nivel como el que el INAI ejecuta «en forma directa». El dictamen del consultor de la misma Subsecretaría, del 4 de septiembre, en el mismo expediente y citando ese informe, dice que el relevamiento se hizo «desde el área de \textbf{DESCENTRALIZADA} del INAI». Dos documentos del mismo organismo, con veinticinco días de diferencia, dan las dos respuestas. El parte provincial describe, para Salta, un \textbf{Equipo Técnico Operativo mixto}, con técnicos de la provincia, del INAI y representantes indígenas, y una resolución del INAI de 2021 confirma que en esta provincia se optó por la ejecución descentralizada. Pero el programa nacional prevé también una \textbf{ejecución centralizada}, a cargo de un equipo central del INAI, que el convenio de 2022 con la Provincia sigue reservando; y en junio de 2013 el INAI y el gobierno provincial todavía estaban acordando la implementación del relevamiento en Salta. Un relevamiento de noviembre de 2011, anterior a esos acuerdos, pudo hacerse de cualquiera de las dos maneras. \textbf{Este libro no lo resuelve: lo resuelve la carpeta técnica o el INAI}, al que el propio informe del 10 de agosto remite «en caso de mayor detalle».

\textbf{Segunda: la carpeta técnica está en manos de la comunidad desde 2016.} No hay que
reconstruirla ni pedirla al archivo nacional: se entregó en un acto público, con funcionarios
provinciales y del INAI presentes. \textbf{La vía más directa para conocer el perímetro
georreferenciado de la zona alta de esta cuenca no es un organismo: es la Asociación
Comunitaria Kondorwaira.}

\textbf{Y tercera: el nombre completo es Asociación Comunitaria Kondorwaira.} El informe de
2023 la nombra «Comunidad Originaria Kolla Kondorwaira»; el parte de 2016, «Asociación
Comunitaria Kondorwaira». Este libro no puede establecer cuál es la denominación registral
exacta, y usa la primera por ser la que consta en el expediente judicial.

\begin{cautela}
Entre el relevamiento de campo --- 20 de noviembre de 2011 --- y la entrega de la carpeta ---
25 de noviembre de 2016 --- pasaron \textbf{cinco años}. Este libro no sabe qué ocurrió en el
medio ni por qué la entrega demoró ese tiempo.

Lo que sí puede señalar es la secuencia: \textbf{el territorio estaba relevado en 2011,
aprobado en 2013 y en manos de la comunidad en 2016}; los fraccionamientos que el capítulo \ref{cap:loteo} documenta se aprobaron entre 2014 y 2021. \textbf{La información existía antes que las
aprobaciones.}
\end{cautela}

\subsection*{Y el territorio tiene número de catastro}

La carpeta técnica no es pública, pero un trabajo académico de 2024 cita su contenido central.

\begin{ficha}{Pérez Domínguez, K.\ M.\ y Belmonte, S., «Arriba están los cerros y abajo la ciudad. Categorizaciones locales y construcción territorial de una comunidad originaria kolla en la provincia de Salta», \textit{Huellas}, UNLPam, 2024}
El área rural de la comunidad es «\textit{la finca Potrero de Castilla (también conocida como finca Las Nieves), identificada con el \textbf{catastro Nº102}, formalmente caracterizado como \textbf{propiedad fiscal de la provincia de Salta} y reclamado como territorio comunitario en la Carpeta Técnica}», con cita de la carpeta como «Archivo Kondorwaira, 2015». Rasgos de la misma forma de vida se observan en las fincas privadas colindantes: \textbf{San Alejo, Torino, Pérez Alsina y Santa Rufina}.

La comunidad \textbf{integra a 31 familias}, se constituyó legalmente en 2009 y obtuvo su carpeta técnica en 2015. Cría ganado menor y mayor con «cambio de puesto» estacional y cultiva papas, maíz, ocas y otros productos andinos. Su sede está en Vaqueros; abajo, sus familias viven en Vaqueros, La Caldera y barrios del norte de la capital, entre ellos \textbf{Juan Manuel de Rosas}.

\textit{DOI 10.19137/huellas-2024-2819; texto completo en \url{https://cerac.unlpam.edu.ar/index.php/huellas/article/download/8217/9710?inline=1}, consultado el 18 de septiembre de 2026. Las autoras son las compiladoras de} Nuestro tiempo redondo.
\end{ficha}

\begin{cautela}
\textbf{Tres consecuencias, y una obliga a corregir.}

\textbf{La pregunta abierta más importante de este libro se vuelve una operación catastral.} Si el núcleo del territorio reclamado es el catastro 102, fiscal, la superposición con los loteos se verifica cruzando esa parcela ---y las cuatro fincas linderas que el artículo nombra--- con las matrículas de los fraccionamientos, sobre la capa catastral de IDESA que este libro ya usó. San Alejo y Santa Rufina, además, son fincas que el libro sigue desde 1909 (capítulos \ref{cap:fincas} y \ref{cap:siglo}). El perímetro georreferenciado sigue en la carpeta; el catastro, no.

\textbf{Dos cifras no coinciden con las de este capítulo.} El parte de 2017 habla de veintisiete familias y el artículo de treinta y una; puede ser crecimiento de la comunidad en siete años. Y el artículo fecha la carpeta en 2015, mientras el parte provincial fecha su entrega en noviembre de 2016: puede ser la fecha de la carpeta frente a la del acto. Este libro consigna las dos y no elige.

\textbf{Y el barrio Juan Manuel de Rosas vuelve a aparecer.} Es el barrio de la capital donde el centro vecinal sostuvo que se adjudicaron lotes a vecinos del departamento (capítulo \ref{cap:tierrafiscal}). Que familias de la comunidad vivan allí no verifica aquella afirmación, pero muestra que el vínculo entre el barrio y el departamento existe por una vía distinta de la que se discutió.
\end{cautela}

\subsection*{El cruce: ningún loteo documentado toca el catastro 102}

Con el número de catastro, la pregunta se pudo contestar sobre la capa «Catastro Urbano y Rural» de la Dirección General de Inmuebles, descargada del geoportal de IDESA en septiembre de 2026, la misma que usan las láminas \ref{fig:planos} y \ref{fig:otbn}. La lámina \ref{fig:kondorwaira} muestra el resultado.

\begin{figure}[p]
\centering
\includegraphics[width=\textwidth,height=0.80\textheight,keepaspectratio]{fig-kondorwaira-catastro.png}
\caption[El territorio reclamado por la comunidad Kondorwaira y los loteos del libro]{\textbf{El territorio reclamado por la comunidad Kondorwaira y los loteos del libro.} En verde, el catastro 102, «Las Nieves», que Pérez Domínguez y Belmonte (2024) identifican como el territorio reclamado en la carpeta técnica; en naranja, las parcelas que lo lindan según la capa catastral; en rojo, los loteos y clubes de campo que documenta este libro; en gris oscuro, las demás parcelas urbanas o subrurales. La trama punteada marca lo que queda fuera de la capa de cuencas del informe técnico de la Ley 8483, es decir, sin categoría de bosque. Elaboración propia con la capa «Catastro Urbano y Rural» de la Dirección General de Inmuebles (geoportal de IDESA, descargada en septiembre de 2026; el organismo la declara orientativa), la cartografía de radios del Censo 2022 (INDEC, corregida por CEUR-CONICET), la capa de cuencas del informe técnico de la Ley 8483 y la hidrografía del IGN. Superficies en proyección equivalente (EPSG:6933). Elaboración propia; dibujo de \texttt{scripts/mapa\_kondorwaira.py} del repositorio \texttt{mapas\_caldera} (apéndice \ref{cap:fuentes}).}\label{fig:kondorwaira}
\end{figure}

\begin{ficha}{Catastro 102 y su entorno, sobre la capa catastral de IDESA}
\begin{itemize}[leftmargin=*,nosep]
\item \textbf{Catastro 102, finca «Las Nieves»}, categoría rural, plano Nº 2: \textbf{22.860 hectáreas} en dos polígonos. Es la parcela más grande del departamento y equivale a cerca de la quinta parte de su superficie.
\item \textbf{Radios censales:} 17.941 hectáreas caen en el radio rural 0101 de Vaqueros ---el 96,5 por ciento de ese radio, donde el censo contó setenta y ocho habitantes--- y 4.067 en el radio rural 0201 de la fracción de La Caldera. Unas 800 hectáreas quedan fuera del contorno censal del departamento, al oeste.
\item \textbf{Parcelas linderas:} San Alejo -- Santa Rufina (catastro 1539, 13.927 hectáreas, plano 120); Lesser (catastro 4416, 10.435 hectáreas, plano 936), y el fraccionamiento Yacones -- Abra de Lesser (catastros 3414 a 3418, plano 615).
\item \textbf{Loteos del libro:} ninguno toca el catastro 102 ni sus parcelas linderas. El más próximo es el remanente rural de El Durazno (catastro 6713), a 6,5 kilómetros; los lotes de El Durazno, el club de campo Finca Lesser, La Milagrosa y El Nogalar están a más de diez. La parcela urbana o subrural más cercana de cualquier clase es una de 2,5 hectáreas (catastro 3419) dentro del fraccionamiento Yacones -- Abra de Lesser, a 3,7 kilómetros.
\end{itemize}
\textit{Distancias medidas entre bordes de parcela en coordenadas UTM 20 S. Los catastros 4366, 3709, 4065, 3616, 1713 y 1905 ya no figuran en la capa: se subdividieron, y sus derivados entran en el cálculo por su número actual cuando el libro lo conoce.}
\end{ficha}

\begin{cautela}
\textbf{Lo que el cruce contesta y lo que no.}

\textbf{Contesta la pregunta tal como estaba formulada, a nivel de catastro, y la respuesta es negativa}: ninguno de los fraccionamientos que este libro documenta se superpone con la parcela que la carpeta técnica reclama, ni con las que la rodean. Los loteos están en el valle; el territorio comunitario, en la cabecera, del otro lado de las fincas grandes.

\textbf{No contesta el perímetro fino.} La carpeta técnica puede reclamar más o menos que la parcela catastral, y el artículo registra que la vida de la comunidad se extiende también a fincas privadas colindantes. La capa, además, es orientativa. La respuesta definitiva sigue en la cartografía de la carpeta.

\textbf{Y deja una pregunta nueva, que mira hacia adelante.} El catastro 1539, San Alejo -- Santa Rufina, lindero al territorio por el este, es la matrícula donde \textsc{agro san alejo s.a.s.}\ fijó su sede en enero de 2026, con un objeto que incluye el turismo rural y los desarrollos inmobiliarios (capítulo \ref{cap:trabajo}). Este libro no sabe si la sociedad es titular o arrendataria de la finca, ni qué piensa hacer en ella. Pero si el próximo fraccionamiento del departamento ocurre ahí, ocurrirá sobre una finca donde, según el artículo, la comunidad también vive y pastorea.
\end{cautela}

\subsection*{Qué contiene una carpeta técnica, y dónde está}

El relevamiento de la Kondorwaira se aprobó por la Resolución INAI Nº 617/2013. Este libro no
tiene esa resolución, pero sí el modelo: las resoluciones de relevamiento del INAI tienen
todas la misma estructura, y una de ellas permite saber exactamente qué documentos existen
detrás de la de 2013 y dónde buscarlos.

\begin{ficha}{Estructura de una resolución de relevamiento del INAI --- Ley 26.160}
El \textbf{artículo 1} da por cumplido el Relevamiento Técnico, Jurídico y Catastral.

El \textbf{artículo 2} «reconoce la ocupación actual, tradicional y pública» de la comunidad
«respecto de la \textbf{superficie georreferenciada, que figura como ANEXO I} y forma parte
integrante de la presente medida».

Y al pie: «El/los \textbf{Anexo/s} que integra/n esta Resolución \textbf{se publican en la
edición web} del Boletín Oficial de la República Argentina».

Los considerandos enumeran las piezas de la \textbf{carpeta técnica}: acta de consentimiento
de la comunidad; \textbf{Cuestionario Socio-Comunitario Indígena}; \textbf{Narrativa} y
\textbf{Croquis}; \textbf{cartografía georreferenciada} con demarcación de los usos internos y
del perímetro; \textbf{Informe Histórico Antropológico}; \textbf{Dictamen Jurídico} con
estrategias de regularización dominial según «la situación registral en que se encuentra el
territorio georreferenciado»; e \textbf{Informe Técnico Final}.
\end{ficha}

Tres consecuencias, y la primera explica por qué el pedido más importante de este libro es difícil de satisfacer.

\textbf{Primera: el perímetro georreferenciado existe, pero probablemente no sea público.} Las resoluciones más recientes de esta clase incorporan la superficie georreferenciada como Anexo I; no se sabe si la de 2013 lo publicó, y la búsqueda de este libro no la encontró (véase la última cautela del capítulo). Además, el convenio de 2022 entre el INAI y el Ministerio de Desarrollo Social de Salta, publicado como anexo del Boletín provincial, declara \textbf{secreta} toda la información del programa de relevamiento, en los términos de la Ley 24.766, y reserva al INAI con exclusividad el derecho de publicarla. \textbf{La vía realista para conocer la cartografía es la propia comunidad.} Con esa cartografía puede establecerse si alguno de los loteos que este
libro documenta se superpone con territorio de ocupación tradicional.

\textbf{Segunda: existe un dictamen jurídico sobre la situación registral del territorio.} La
carpeta técnica incluye, por diseño, un análisis de quién figura como titular registral de
las tierras relevadas y qué estrategias caben para regularizarlas. \textbf{Es, para la zona
alta de esta cuenca, el informe de dominio que el capítulo \ref{cap:loteo} reclama para las matrículas del
loteo} --- hecho por el Estado nacional y diez años antes.

\textbf{Y tercera: el reconocimiento no es declarativo de propiedad.} El artículo 2 reconoce
la \textit{ocupación}, no titula. La Ley 26.160 declara la emergencia y suspende desalojos;
la propiedad comunitaria requiere una ley que el Congreso no dictó. \textbf{Lo que existe
sobre la zona alta de esta cuenca es un reconocimiento estatal de ocupación, no un título}, y
este libro no debe leerlo como si lo fuera.

\subsection*{Tres precisiones del expediente}

\textbf{El relevamiento tiene fecha exacta: 20 de noviembre de 2011.} El informe de 2023 lo califica de \textbf{«centralizado»}, y una primera lectura de este libro entendió que lo había hecho el equipo del INAI en forma directa. La carpeta técnica se aprobó por Resolución INAI Nº 617/2013.

El informe del 10 de agosto usa la palabra para la ejecución ---«fue realizado por el equipo técnico del INAI»---, pero el dictamen del 4 de septiembre, en el mismo expediente, lo atribuye al «área de DESCENTRALIZADA del INAI». Como se dijo más arriba, este libro no puede establecer cuál de los dos tiene razón.

\textbf{La segunda comunidad tramita su personería desde 2018 y el trámite está en el
Instituto Provincial de Pueblos Indígenas.} Cinco años después de iniciado, el expediente
54-142017/2018-0 seguía ahí en agosto de 2023.

\textbf{Y la tercera es la que más dice, aunque no hable de comunidades sino del Estado.}

Quien pidió el informe fue la \textbf{abogada de la Fiscalía de Estado de la Provincia}, en
representación de la propia Provincia demandada, y lo pidió «en carácter urgente»,
acompañando «mapa imagen satelital del área respectiva». El pedido dice, textualmente, que se
informe «si en la zona \textbf{se encuentra relevado} la existencia de comunidades
aborígenes».

Es decir: \textbf{en agosto de 2023, para saber si había comunidades en la cuenca que estaba
obligada a ordenar, la Provincia tuvo que preguntarse a sí misma}, girando el expediente de
un ministerio a otro y adjuntando una imagen satelital para que le identificaran el área.

El dato existía desde 2011 y estaba aprobado por resolución desde 2013. \textbf{No fue
necesario producirlo: fue necesario encontrarlo.} Y lo encontró, en una semana, una dirección
de otro ministerio de la misma provincia.

Este libro sostiene en el capítulo \ref{cap:opacidad} que la opacidad de este dispositivo rara vez consiste
en negar un documento. Aquí hay una variante que no había registrado: \textbf{un organismo del
Estado que no sabe lo que otro organismo del mismo Estado tiene documentado hace doce años.}

\begin{cautela}
Lo que este libro puede afirmar ahora y lo que sigue sin poder afirmar.

\textbf{Puede afirmar} que las comunidades existen, que están registradas, que su territorio
está relevado y aprobado, y que fueron convocadas a los talleres del Plan.

\textbf{No puede afirmar} cuántas personas las integran, qué superficie abarca el relevamiento
de 2011 ---la prensa oficial habla de veintisiete familias--- ni si el Plan aprobado recoge lo que las comunidades plantearon en la consulta, cuya acta transcribe el capítulo \ref{cap:amparo}.
Sobre la superposición con los loteos, el cruce catastral de este mismo capítulo da una respuesta negativa a nivel de parcela; el perímetro fino sigue en la carpeta técnica, y la finca lindera de San Alejo es, desde 2026, \textbf{la pregunta abierta más importante de este trabajo}.

El capítulo \ref{cap:infraestructura}, más adelante, enuncia como una de sus cuatro condiciones de refutación que aparecieran comunidades con ocupación anterior al archivo. \textbf{Aparecieron, aunque con un reconocimiento estatal de ocupación y no con un título.} Lo que la
tercera tesis decía --- que el mercado encontró un dispositivo armado --- hay que leerlo ahora
sabiendo que sobre ese mismo territorio había, además, una ocupación anterior a todo el
archivo que este libro recorrió.
\end{cautela}

\subsection*{Quince años de reclamo}

La cobertura consigna antecedentes de reclamo vecinal en La Caldera por la protección de los ríos y su entorno desde 2010: protestas, notas y denuncias ante la entonces Secretaría de Ambiente provincial y ante el municipio, y finalmente amparos colectivos.

La cronología del apéndice \ref{cap:cronologia} permite reconstruir la progresión: reclamos desde 2010, advertencia administrativa en 2016, la gacetilla de octubre de 2018 que este capítulo transcribe, denuncia pública en 2019, acción judicial en 2019, sentencia favorable en 2021, y ejecución litigada entre 2022 y 2026. Es, en escala municipal, el recorrido completo de los grados de disidencia que el libro anterior describió.

\section{Lo que resultó eficaz}

\textbf{El amparo colectivo ambiental.} Obtuvo lo que once años de gestión administrativa no habían obtenido: una condena a confeccionar un plan ---cuyo alcance sobre la ejecución de obras discuten las partes (capítulo \ref{cap:amparo})---, con criterio técnico definido por el tribunal y seguimiento judicial de la ejecución.

\textbf{La producción de conocimiento técnico independiente.} Una vecina con formación en biología produjo el registro público más citado del proceso de desmonte, sosteniendo que el negocio inmobiliario operaba sin control y con consecuencias de gravedad impredecible. Ese registro fue levantado por prensa nacional y funcionó como pieza de legitimación del reclamo.

\textbf{La peritación profesional colegiada.} El Consejo Profesional de Agrimensores, Ingenieros y Profesiones Afines actuó como veedor en la etapa de ejecución y recibió reconocimiento expreso del tribunal.

\textbf{El perfil de quienes litigaron.} El capítulo \ref{cap:poblacion} ubica al loteo El Durazno en el radio 0208 o en el 0209, y cualquiera de los dos está entre los de mayor formación del municipio: 21 y 24 personas con estudios universitarios por cada cien adultos, por encima de la capital, que tiene 16,7; en el 0209, más de la mitad de los adultos cursó un terciario o una carrera universitaria. El conjunto de la fracción queda en tres de cada diez \textit{contando el terciario} ---y en 13,6 si se cuentan sólo los universitarios---, y el radio más poblado del municipio, en 5 universitarios por cada cien. De modo que el nivel educativo del entorno de El Durazno es alto, sea cual fuere el radio. Con las reservas de aquel capítulo ---radios de 125 y 150 habitantes, propietarios de fin de semana que no se cuentan allí y barrios ubicados por inferencia---, los loteos que sostuvieron siete años de amparo parecen ser los que tenían, en proporción, más herramientas para hacerlo. No todos: El Nogalar, de origen municipal, también es parte actora, y este libro no sabe quiénes lo habitan.

\begin{cautela}
El rasgo compartido es la independencia estructural respecto del aparato local. Ni el amparista, ni la científica, ni el colegio profesional dependen del municipio para su reproducción material o simbólica. Es el mismo principio que el libro anterior formuló al cerrar su inventario de disidencias no absorbidas: lo que el dispositivo no puede absorber es aquello que no necesita su reconocimiento para existir.

El dato del censo agrega un matiz que conviene no omitir: esa independencia también tiene una base social. Que los más formados resistan mejor no le quita mérito a su resistencia, pero sí advierte sobre los ausentes. Los barrios del pueblo con menos formación sufrieron el mismo río y no aparecen en ningún expediente.
\end{cautela}

\section{Los límites}

Conviene no escribir una épica. El amparo de Vaqueros contra el loteo eclesial fue rechazado en 2019 y el emprendimiento avanzó. El amparo de La Caldera fue ganado en 2021 y las inundaciones volvieron en enero de 2026, semanas antes de que el plan fuera finalmente aprobado. Los reclamos por agua potable en la zona alta de Vaqueros llevan más de quince años sin resolución. El derrumbe de ladera sobre la ruta y las denuncias del centro vecinal de La Calderilla muestran que el frente de desmonte no se detuvo. Y el capítulo \ref{cap:amparo} documentó que la condena reapareció, cuatro años después, como propuesta de gestión del propio condenado.

\begin{cautela}
Una sentencia favorable obtenida contra un municipio sin capacidad técnica ni recursos para cumplirla produce una obligación, no un resultado. Ése es el límite estructural de la vía judicial en esta escala, y el libro no debe disimularlo.

De ahí se sigue una consecuencia práctica para quien litigue: el objeto útil de un amparo en este departamento no es sólo la condena al municipio, sino la condena a la Provincia a proveerle los medios. La Fiscalía de Estado ya estableció, en el propio expediente, que la responsabilidad de las obras es municipal porque el municipio autoriza. Falta que alguien establezca con qué debe hacerlas.
\end{cautela}

\section{Los ausentes}

Ninguno de los siguientes sujetos aparece en la cobertura periodística relevada, ni como parte en los expedientes identificados, ni en la literatura académica consultada: la población rural preexistente; los trabajadores de la construcción de los propios loteos; los ocupantes sin título; y las familias con posesión veinteañal no comunitaria. Las comunidades originarias estuvieron en esta lista hasta que el expediente de Fiscalía de Estado permitió sacarlas de ella.

\begin{cautela}
Su ausencia del corpus no prueba su ausencia del territorio. El libro anterior documentó, para otros departamentos, exactamente ese mecanismo: la invisibilidad documental como primer paso de la invisibilidad jurídica.

La diferencia entre ``verifiqué y no hay'' y ``no encontré'' es todo el método de este libro. Para las comunidades originarias ese paso ya se dio: existen y están registradas, como documenta el comienzo de este capítulo. Para los otros ausentes de la lista, todavía no.

Hay, sin embargo, dos rastros que conviene retener.

El primero es cuantitativo y el capítulo \ref{cap:tierra} lo documenta: el Censo Nacional Agropecuario 2018 registra en el departamento sesenta y una explotaciones agropecuarias sin límites definidos, el 46,6 por ciento del total, contra el 32,8 provincial. Esa categoría no equivale a tenencia irregular, pero es el rastro estadístico más próximo a la población que este apartado nombra como ausente.

El segundo es cartográfico y es más directo. El informe técnico de la Ley 8483 establece que el criterio 11 de la Ley 7543 --- ``valor y uso dado por comunidades indígenas y campesinas a áreas boscosas o colindantes'' --- funciona como \textbf{capa de control de gestión}, y produce un mapa específico, identificado como Mapa 11 en el Anexo IX del informe. Cada proyecto de manejo o cambio de uso del suelo debe someterse a un protocolo de consulta previa, libre e informada, y no puede aprobarse sin el visto bueno de las Comisiones de Evaluación conformadas para aplicarlo.

Existe, por lo tanto, un mapa oficial y público que registra dónde hay comunidades indígenas y campesinas en la provincia de Salta. Consultado para este libro, el mapa muestra el departamento bajo la mancha de uso potencial indígena, pero a su escala no permite saber si el territorio relevado de la Kondorwaira figura en él ni con qué alcance.

Hay, sin embargo, un dato del propio proceso que conviene poner junto a esa pregunta. El Consejo Asesor que produjo el ordenamiento clasificó a sus cincuenta y tres integrantes en siete tipos, y el séptimo --- \textit{representantes de pueblos originarios} --- registra \textbf{cero}, según el informe técnico que integra el Anexo I de la Ley 8483. En el taller del 25 de octubre de 2023, en cambio, se acreditaron diecisiete representantes de pueblos originarios, casi uno de cada cinco participantes: la ausencia es de instituciones indígenas en el Consejo, no de personas en el taller.

Un libro que se pregunta por los ausentes de su territorio encuentra, en el instrumento provincial que debía identificarlos, la misma ausencia formalizada en una tabla.
\end{cautela}

\section{Una asociación anterior al conflicto}

Entre las convocatorias a asamblea que publica el Boletín Oficial aparece una entidad que este libro no había registrado y que importa por su fecha.

\begin{ficha}{Convocatoria a asamblea ordinaria --- B.O.\ Nº 19.430, 19 de noviembre de 2014}
La \textbf{``Asociación de Vecinos de los Ríos Castellanos y Lesser''} convoca a sus socios a asamblea ordinaria para el 23 de diciembre de 2014, en calle M.\ Julia Solá 309 de \textbf{Lesser}, para considerar memoria, balance, inventario e informe de la comisión revisora de cuentas, y \textbf{elegir todos los miembros} de la comisión directiva y de la revisora de cuentas.

Firman por la comisión directiva \textbf{Inés Ortíz de Cárdenas} y \textbf{José Armando Caro Figueroa}.
\end{ficha}

\begin{cautela}
Tres puntos que este dato aporta al capítulo.

\textbf{La primera es cronológica.} En noviembre de 2014 existía en el departamento una asociación civil con personería, comisión directiva, órgano de fiscalización y balance anual, \textbf{organizada explícitamente alrededor de dos ríos}. Cuatro años antes del primer aluvión y más de cuatro antes del amparo colectivo. La organización vecinal en torno al agua no nació del desastre: lo precedía.

\textbf{La segunda es de escala.} Los ríos Castellanos y Lesser corresponden a parajes del departamento que este libro apenas menciona, y que no figuran en ninguno de los conflictos documentados. Que allí hubiera vecinos asociados y renovando autoridades sugiere que el tejido asociativo del corredor es más denso de lo que la reconstrucción por conflictos permite ver. Un libro que sigue litigios ve sólo donde hubo litigio.

\textbf{La tercera es un nombre.} Uno de los firmantes, \textbf{José Armando Caro Figueroa}, es homónimo de quien fuera Ministro de Trabajo de la Nación entre 1993 y 1997, salteño de origen. Este libro \textbf{no afirma que sean la misma persona}: comparte nombre completo, lo que es indicio y no prueba, y la coincidencia se resuelve con un padrón, no con una lectura. Se consigna porque, de confirmarse, indicaría que la asociación de vecinos de esos ríos contaba con integrantes de peso institucional considerable, y eso cambiaría la lectura de su capacidad de acción.
\end{cautela}

Y el relevamiento de los boletines de 2015 confirma la densidad. En marzo aparecen convocando a asamblea la \textbf{Liga Municipal de Fútbol San Lorenzo -- La Caldera}, con sede en avenida San Martín de Vaqueros, tratando balances de los períodos 2010 a 2013 y la incorporación definitiva de clubes; y el \textbf{Club Social y Deportivo de Caza y Pesca La Caldera}, que sesiona en el local del Centro de Jubilados Cristo Monumental, en calle Las Margaritas y Madreselva. Con la Asociación de Vecinos de los Ríos Castellanos y Lesser, el Centro de Jubilados de Vaqueros y la Asociación Civil CESPS --- que en enero de 2015 renovaba comisión directiva en General Güemes 250 de Vaqueros ---, son al menos cinco entidades con personería y vida asamblearia en un departamento que hoy tiene doce mil habitantes.

Dos observaciones. \textbf{La primera:} el nombre de la liga de fútbol --- San Lorenzo y La Caldera juntos, con sede en Vaqueros --- vuelve a mostrar que el corredor funciona como una unidad social que ninguna jurisdicción administra entera. Lo que el capítulo \ref{cap:redes} documentó para el transporte y para el río, la liga de fútbol lo resolvió sola en 2015.

\textbf{La segunda:} varias de esas asambleas tratan balances atrasados de tres y cuatro ejercicios. Es un rasgo común de las entidades pequeñas y no indica irregularidad, pero sí que la vida asociativa del departamento es real y a la vez precaria: existe, se sostiene y le cuesta cumplir sus propias formas.

\begin{cautela}
Dos entidades más, y una de ellas obliga a mirar el capítulo \ref{cap:tierra} de otra manera.

En junio de 2015 convoca a asamblea la \textbf{Asociación Civil de Pequeños Productores Agropecuarios del Departamento La Caldera}, con sede en calle Las Moreras 587 del municipio de Vaqueros, para tratar memoria y balances de 2013 y 2014 y renovar autoridades. Firma como secretario el Dr.\ Diego A.\ M.\ Ragone.

\textbf{Existían pequeños productores agropecuarios organizados en el departamento en 2015}, y su asociación siguió convocando asambleas al menos hasta 2024 (capítulo \ref{cap:ausencias}). El capítulo \ref{cap:tierra} documenta que entre 2002 y 2018 el departamento perdió el 88,4 por ciento de su superficie agrícola, mientras la provincia ganaba terreno. Ese desplome no ocurrió sobre un vacío: ocurrió sobre productores que tenían asociación con personería, órgano de fiscalización y cuota social.

Este libro no sabe cuántos eran, qué producían ni qué fue de ellos. Pero la existencia de la entidad cambia el modo de formular la pregunta del capítulo \ref{cap:tierra}. Ya no es ``qué pasó con las hectáreas'': es \textbf{qué pasó con la gente que las trabajaba}, y hay un registro donde empezar a averiguarlo.

La otra entidad es el \textbf{Club de Regatas Güemes}, cuya sede está \textbf{en el Dique Campo Alegre}. Es decir, sobre el espejo de agua que la Ley 4486 hizo posible expropiando cinco propiedades en 1972, y en cuyo entorno el municipio proyectaba en 2026 un ``nuevo circuito turístico''. Un club náutico con comisión directiva y balances, instalado en el embalse desde hace décadas, es un actor del área que este libro no había registrado como tal: el capítulo \ref{cap:tierrafiscal} documenta el permiso precario de 1980 a otro club náutico del perilago, el Club de Regatas Campo Alegre, y no establece si son la misma entidad con otro nombre.
\end{cautela}

\section{Un Estado que sí llega: la mediación comunitaria}

\begin{ficha}{Decretos Nº 475 y 481 del 9 de febrero de 2015 --- B.O.\ Nº 19.485}
El Ministerio de Justicia prorroga por doce meses, a partir del 1 de enero de 2015, los contratos de locación de servicios de \textbf{dos mediadoras comunitarias} que se desempeñan en el \textbf{Centro de Mediación Comunitaria de La Caldera}: \textbf{Mariángeles Giraldez}, cuyo contrato se aprobó por Decreto 3722/09, y \textbf{Flavia del Valle Tabarcache}, por Decreto 3871/10.

Ambos contratos dependen de la Secretaría de Métodos Alternativos de Resolución de Conflictos.
\end{ficha}

\begin{cautela}
Es un dato pequeño y corrige un sesgo de este libro.

El departamento tiene, desde 2009 y 2010, un \textbf{Centro de Mediación Comunitaria} con dos profesionales contratadas por la Provincia y renovadas año tras año. Existe también, según los mismos boletines, un centro equivalente en Vaqueros. Es decir: hay una política pública provincial de resolución de conflictos instalada en el territorio, sostenida durante más de un lustro.

Este libro se construyó siguiendo litigios, desmontes y expedientes de aprobación, y esa elección produce una imagen del Estado que llega tarde o no llega. La mediación comunitaria es un caso donde llegó, se quedó y se renovó. No hace falta que eso equilibre el balance --- no lo equilibra ---, pero sí que quede registrado, porque un libro que sólo mira donde hubo conflicto termina describiendo un Estado que sólo existe en el conflicto.

Lo que no sabemos es qué se medió. Un centro de mediación comunitaria en un departamento donde se disputan tierras, aguas y loteos pudo intervenir en cualquiera de esos frentes, o en ninguno. Sus registros lo dirían.
\end{cautela}

\section{Áridos sobre tierra fiscal}

\begin{ficha}{Edicto de mina --- Expte.\ Nº 18.219, publicado en enero y febrero de 2014}
El Juzgado de Minas y en lo Comercial de Registro, a cargo del Dr.\ Daniel Enrique Marchetti, hace saber que \textbf{Transporte Hnos.\ Mendoza S.R.L.} solicitó la publicación de la \textbf{renovación de concesión} de la cantera de áridos denominada \textbf{``Macarena''}, ubicada en el departamento La Caldera, \textbf{lugar río Wierna}.

Superficie libre: \textbf{10 hectáreas}. Coordenadas Gauss Krüger, sistema Posgar 94 y Campo Inchauspe/69.

\textbf{``Los terrenos afectados son de propiedad Fiscal.''}
\end{ficha}

\begin{cautela}
Es el primer edicto minero del departamento que este relevamiento encontró para el período 2014--2015, y reúne tres hilos que este libro llevaba separados.

\textbf{El primero:} hay extracción de áridos concesionada sobre el río Wierna, uno de los cursos que --- según el testimonio vecinal recogido en el capítulo \ref{cap:expropiacion} --- rodean Vaqueros y cuya agua los vecinos no reciben.

\textbf{El segundo:} el artículo 171 del Código de Aguas exige que, \textbf{antes} de otorgar permisos de extracción de áridos, la Autoridad de Aplicación evalúe si afectarán desfavorablemente las riberas o el flujo de las aguas, y en su caso desaconseje el otorgamiento o exija obras. Una \textbf{renovación} de concesión es una nueva instancia de otorgamiento. Si esa evaluación se hizo, está en el expediente.

\textbf{Y el tercero, que es el que más pesa:} los terrenos son de \textbf{propiedad fiscal}. Diez hectáreas de tierra pública en un departamento donde, según el capítulo \ref{cap:tierrafiscal}, el circuito de adjudicación de tierra fiscal a familias tarda décadas y exige nacionalidad, cuatro años de residencia y no poseer inmuebles en ningún lugar del mundo. La misma tierra pública se concesiona para extraer áridos con un edicto de tres publicaciones.

Este libro no compara los dos regímenes para reprochar el segundo: la concesión minera y la adjudicación social son institutos distintos, con finalidades distintas, y ninguno es ilegítimo. Los compara porque \textbf{ésa es exactamente su tesis}: sobre el mismo suelo público operan dos circuitos con intensidades regulatorias inversamente proporcionales a la capacidad de sus beneficiarios.
\end{cautela}

En febrero y marzo de 2014, casi en simultáneo con el edicto de «Macarena», el mismo juzgado publicó el edicto de renovación de la cantera de áridos \textbf{``El Vaquero''}, expediente 19.199/08, solicitada por \textbf{INCOVI S.R.L.}, ubicada en el departamento La Caldera, lugar Vaqueros, con una superficie de \textbf{nueve hectáreas y 347 metros cuadrados}. También allí: ``los terrenos afectados son de propiedad Fiscal''.

Y hay una tercera, sobre el mismo curso aunque del otro lado del límite: en abril de 2014 se publicó la renovación de la cantera \textbf{``El Triunfo II''}, expediente 17.422, solicitada por la \textbf{Cooperativa de Provisión de Áridos La Mesa Redonda Ltda.}, ubicada ``en el departamento Capital, río La Caldera''. El yacimiento está en la Capital; el río es el mismo.

Tres concesiones de áridos sobre la misma cuenca en un solo trimestre, dos de ellas sobre \textbf{tierra fiscal} dentro del departamento y sumando casi veinte hectáreas. Las dos primeras por renovación, con antecedentes que se remontan al menos a 2008.

Y la tercera agrega el dato jurisdiccional que el capítulo \ref{cap:vaqueros} ya había encontrado por otra vía: el río La Caldera y su continuación atraviesan dos departamentos, y la extracción de áridos se autoriza a ambos lados sin que ningún organismo tenga competencia sobre el conjunto del cauce.

El contraste con el circuito de adjudicación social: son las mismas hectáreas fiscales, la misma provincia y el mismo año.

Y no es una sola cantera. Para noviembre de 2015 el relevamiento suma \textbf{seis concesiones de áridos} en la cuenca ---una de ellas en jurisdicción de la Capital---: ``Macarena'' (Transporte Hnos.\ Mendoza, río Wierna, 10 ha fiscales), ``El Vaquero'' (INCOVI, Vaqueros, 9 ha fiscales), ``El Triunfo II'' (Cooperativa La Mesa Redonda, río La Caldera, jurisdicción Capital), \textbf{``Los Yacones''} (INCOVI, cauce del río Wierna, \textbf{21 ha 6.374 m² de terrenos fiscales}) \textbf{``Don Roque I''} (M.E.I.\ Obras y Servicios S.R.L., río Vaqueros) y \textbf{``Joaquín''} (René Antonio Barranco Barrera, río Las Nieves).

Con ``Los Yacones'', \textbf{INCOVI S.R.L.\ acumula dos canteras sobre tierra fiscal} --- más de treinta hectáreas entre ambas --- y es la cuarta aparición de esa firma en este libro, después de la cantera de Vaqueros, la oferta en la obra de protección del acueducto y la contratación directa para su reparación. Se consigna, como en los casos anteriores, sin atribuir irregularidad: una empresa de obra vial e hídrica que explota canteras de árido hace lo propio de su rubro.

\subsection*{Y hay minería propiamente dicha}

Las canteras de áridos no agotan la actividad extractiva del departamento.

\begin{ficha}{Mina ``Cristina M-1'' --- Expte.\ 22.295, publicado en abril de 2015}
El Juez de Minas hace saber que \textbf{Roberto Adrián Scasso, Julio Argentino Monteros, Facundo Martín Monteros Moya y Federico Joaquín Monteros Moya} manifestaron el descubrimiento de un \textbf{yacimiento de fosfatos} en el departamento La Caldera, distrito Mojotoro, \textbf{Sierra de Mojotoro}.

\textbf{Superficie concedida: 1.256 hectáreas 3.712 m².}

``Los terrenos son \textbf{privados}'': la matrícula 1.902 pertenece a \textbf{Jorge Raúl Sanz}; las matrículas \textbf{1.981, 1.982, 1.983, 1.984 y 1.985 pertenecen a TURPINAR S.A.}; y la matrícula 1.862 se encuentra sujeta a inscripción registral.

Plazo para que quienes tuvieren derechos los hagan valer: sesenta días.
\end{ficha}

Dos cosas, y la segunda excede a este capítulo.

\textbf{La primera: la escala.} Mil doscientas cincuenta y seis hectáreas es una superficie mayor que la de varios de los loteos y clubes de campo que este libro documenta juntos. Y no se trata de áridos de río sino de un yacimiento de \textbf{fosfatos}, un mineral de segunda categoría con destino industrial. La Sierra de Mojotoro, sobre el límite oriental del departamento, tiene entonces valor minero declarado.

\textbf{La segunda: aparece un propietario corporativo.} El edicto identifica a \textbf{TURPINAR S.A.} como titular de \textbf{cinco matrículas contiguas} --- 1.981 a 1.985 --- en la sierra. Es la primera vez que este relevamiento encuentra a una sociedad anónima concentrando esa cantidad de inmuebles en el departamento, y aparece por vía indirecta: no en un aviso societario ni en una aprobación de loteo, sino en la citación de un expediente minero de tercero.

Es exactamente el tipo de dato que el capítulo \ref{cap:opacidad} explica que no se encuentra: la propiedad concentrada no se publica, se deduce de actos que la mencionan de paso. Quiénes integran TURPINAR S.A., desde cuándo posee esas matrículas y qué superficie suman es materia de la Dirección General de Inmuebles y del legajo societario.

\begin{cautela}
Un resultado negativo de búsqueda, que no contradice lo anterior y conviene explicar. La búsqueda sistemática en buscadores públicos de resoluciones del Instituto Nacional de Asuntos Indígenas que den por cumplido el relevamiento técnico, jurídico y catastral de la Ley 26.160 en la provincia de Salta no arrojó ninguna correspondiente al departamento La Caldera. Las localizadas corresponden a Molinos, Cachi, La Poma, San Carlos, General San Martín y otros departamentos. En Salta esos relevamientos se ejecutaron, por regla general, mediante Equipo Técnico Operativo convenido entre el INAI y el ministerio provincial. La Resolución INAI 617/2013, que aprobó el de la Kondorwaira, se conoce por el informe del Equipo Técnico de 2023 y no apareció en esa búsqueda: es un ejemplo más de por qué este libro no convierte una ausencia en un buscador en una ausencia en el registro.

Lo que puede afirmarse es que no se localizó acto administrativo publicado de relevamiento en el departamento.

Hay además una fuente que contesta la pregunta por vía positiva y que este libro ya sabe procesar: el Censo 2022 publica, por departamento, la población que se reconoce indígena o descendiente de pueblos indígenas u originarios, con desagregación por pueblo a nivel provincial. Es el mismo camino de Redatam que el capítulo \ref{cap:poblacion} empleó para la vacancia habitacional.
\end{cautela}

